% Fichier complété par tools/complete_missing_corrections.py.
% Source : STAT/STAT-02-Frottement/533_MAM_Frottement_Cylindre/533_MAM_Frottement_Cylindre.tex
% Statut : rédigé (3 question(s)).
\begin{corrigebox}[Corrigé rédigé]
\CorrigeQuestion{1}
Au contact cylindre/plan, la réaction se décompose en une normale $N$ et
            une tangente $T$. La loi de Coulomb impose $|T|\le fN$ en adhérence et
            $T=-fN\,\vec t_v$ au glissement. Le torseur local est donc
            $\{\mathcal T\}=\{T\vec t+N\vec n;\vec0\}_I$.
\CorrigeQuestion{2}
À la limite du glissement, $|T|=fN$ et l'angle de frottement vérifie
            $\tan\varphi=f$. L'équilibre global fournit ensuite $N$ et $T$ par projection
            des forces et moments.
\CorrigeQuestion{3}
Le sens de $T$ est opposé à la vitesse de glissement probable. La
            vérification finale consiste à contrôler que la solution d'adhérence respecte
            $|T|<fN$; sinon on reprend avec l'égalité de glissement.
\end{corrigebox}
